% Drafting owner: Ettore. Keep scientific content student-authored.
\subsection{Background, Rationale, and Related Work}

Seabed mapping can be treated as a coverage problem: a vehicle must plan a path so that the relevant area, or volume, is observed by its sensors. Coverage path planning is a central problem in robotics and includes underwater applications such as image mosaicing and inspection \cite{galceran2013coverage}. In marine settings, useful coverage also depends on sensing footprint, vehicle motion, positioning uncertainty, environmental disturbances, and information exchange during the mission.

Using multiple vehicles can improve survey efficiency because the mission area can be divided across the team. However, marine multi-robot systems introduce coordination problems that do not appear in the same way for a single vehicle. Prior work identifies communication, task division, task allocation, and coordination as closely connected challenges, especially because marine communication links can be limited or intermittent \cite{martorell2024coordination}. If a team replans online, then stale shared information may affect which regions are assigned, revisited, or left for later.

The required simulation fidelity is an open design choice. Surface vehicles can often be represented on a 2D plane, while underwater vehicles may require depth or altitude decisions and more constrained acoustic communication. Recent AUV trajectory-planning literature distinguishes 2D, 2.5D bathymetric, and full 3D representations, with increasing realism and modelling cost \cite{almuzaini2026trajectory}. Underwater acoustic networking also faces limited bandwidth, long and variable delays, Doppler effects, multipath, and environmental dependence \cite{alshawabka2026underwater}. For this project, a controlled model is therefore a suitable starting point because it can isolate the coordination effect before adding realism.
